Attraverso lo studio dei raggi cosmici utilizzando emulsioni nucleari in mongolfiera, nel 1936 fu scoperta l'esistenza di una particella che, in presenza di un campo magnetico deviava in maniera differente rispetto all'elettrone. Questa particella è il \textbf{muone}. L'origine dei muoni osservati sulla superficie terrestre si può far rilsalire al decadimento di un'altra particella: il pione. In questa relazione si vuole presentare la misura della vita media del muone attraverso misure di coincidenza rese possibili grazie a delle proprietà fondamentali del muone e alla presenza di 3 scintillatori organici. 